\documentclass[12pt]{article}
\usepackage[utf8]{inputenc}
\usepackage[T1]{fontenc}
\usepackage[margin=1in]{geometry}
\usepackage{setspace}
\usepackage{xcolor}
\usepackage[colorlinks=true, urlcolor=blue!60!black, linkcolor=blue!60!black]{hyperref}
\setlength{\parindent}{2em}
\setlength{\parskip}{1.25em}
\renewcommand{\baselinestretch}{1.0}
\begin{document}
\begin{center}
{\LARGE \textit{Kelo v. City of New London}, 545 U.S. 469 (2005)\par}
\vspace{0.8em}
{\large Case Brief --- Four-Page Edition\par}
\end{center}
\vspace{1em}

\section{Memory Jogger}

Upholding 5--4 New London's condemnation of unblighted homes under a comprehensive plan for economic development, the Court held that promoting economic development is a ``public purpose'' satisfying the Takings Clause's public-use requirement under deferential review, while leaving States free to impose stricter limits (\href{https://supreme.justia.com/cases/federal/us/545/469/}{Justia}).

\section{Detailed Case Facts}

After decades of decline, including the 1996 closure of the Naval Undersea Warfare Center and its 1,500 jobs, New London, Connecticut, had unemployment nearly double the state average. City officials reactivated the New London Development Corporation (NLDC), a private nonprofit, to revitalize the 90-acre Fort Trumbull peninsula. After Pfizer announced a \$300 million research facility next door in 1998, the NLDC drafted a seven-parcel plan with a waterfront hotel, restaurants, a riverwalk, a marina, 80 new residences, a Coast Guard Museum site, at least 90,000 square feet of research office space (Parcel~3), and a 2.4-acre ``park support'' or marina-support tract (Parcel~4A).

The city council approved the plan in 2000 under Connecticut's municipal development statute and authorized condemnation. The NLDC bought most of the 115 properties from willing sellers but condemned 15 well-maintained, unblighted properties owned by nine holdouts, including Susette Kelo, who had renovated her waterfront home, and Wilhelmina Dery, who had lived in her home since her birth there in 1918.

\section{Procedural History}

The owners sued in New London Superior Court under the Public Use Clause. After a seven-day trial, the court enjoined the Parcel~4A takings as too vague but upheld those on Parcel~3. The Connecticut Supreme Court, 4--3, upheld all the condemnations under \textit{Berman v.\ Parker} and \textit{Hawaii Housing Authority v.\ Midkiff}. The U.S. Supreme Court affirmed on June 23, 2005.

\section{Judicial Votes}

Affirmed, \textbf{5--4}.
\begin{itemize}
\item \textbf{Majority:} Justice Stevens, joined by Justices Kennedy, Souter, Ginsburg, and Breyer.
\item \textbf{Concurring:} Justice Kennedy.
\item \textbf{Dissenting:} Justice O'Connor, joined by Chief Justice Rehnquist and Justices Scalia and Thomas; Justice Thomas separately.
\end{itemize}

\section{Holding}

Condemning private property for transfer to private parties under a carefully formulated, comprehensive economic-development plan serves a valid ``public purpose'' and satisfies the ``public use'' requirement under deferential review, even where the properties are not blighted and will not be opened to the public. States remain free to impose stricter limits under state law.

\section{Analysis of Opinions}

\textbf{Majority (Stevens).} Government may not take property from $A$ solely to transfer it to $B$, or under the ``mere pretext'' of a public purpose, but may transfer it to a common carrier. The Court had ``long ago rejected any literal requirement that condemned property be put into use for the general public,'' adopting ``public purpose.'' Under \textit{Berman}, \textit{Midkiff}, and \textit{Monsanto}, comprehensive plans are judged as a whole, by purpose rather than mechanics. Economic development is a traditional function with no principled distinction from other purposes; the deliberate, state-authorized plan, whose private beneficiaries were mostly unknown, showed no favoritism. Courts need not demand proof to a ``reasonable certainty'' of benefits, and States may add restrictions.

\textbf{Kennedy (concurring).} Rational-basis review still bars takings ``intended to favor a particular private party, with only incidental or pretextual public benefits,'' as in \textit{Cleburne} and \textit{Moreno}. There was no pretext here, but ``a more stringent standard of review'' might apply to suspicious transfers.

\textbf{O'Connor (dissenting).} Equating public use with incidental economic benefit would ``wash out any distinction between private and public use of property.'' \textit{Berman} and \textit{Midkiff} involved harmful precondemnation uses, and Kennedy's pretext test would be failed only by a ``stupid staffer.''

\textbf{Thomas (dissenting).} \textit{Berman} and \textit{Midkiff} should be overruled: government may take only property it will own or the public may legally use. Urban renewal (``Negro removal'') shows deferential review exposes powerless and minority communities to displacement.

\section{Examples: Future Applications}

\begin{enumerate}
\item \textbf{Same-Side: Regional Industrial Park.} After its largest factory closes, a county adopts, after studies and hearings, a 200-acre industrial-park plan with tenants chosen later and condemns holdout farmhouses. The plan is judged as a whole and is valid.
\item \textbf{Same-Side: Common-Carrier Rail Line.} A regulated railroad open to all shippers condemns farmland for a freight line. Transfers to common carriers are settled public use.
\item \textbf{Opposite-Side: One-Parcel Transfer.} Without any plan, a prosperous suburb condemns a profitable nursery two weeks after its owner refused MegaTech's offer, for transfer to MegaTech. Under the ``mere pretext'' bar and Kennedy's concurrence, the taking is likely invalid.
\item \textbf{Opposite-Side: State Constitutional Bar.} A State's constitution forbids takings for private economic development. \textit{Kelo} is only a floor, so the state bar controls.
\item \textbf{Fence-Sitter: Arena for a Named Team.} A city adopts a 20-acre plan, after hearings, built around a team's new arena. The planning process supports public purpose, but a single identified beneficiary invites Kennedy's heightened scrutiny.
\end{enumerate}

\section{Critique}

\textbf{Compensation does not make owners whole.} Stevens assumed that ``just compensation'' mitigates hardship, yet fair-market value omits subjective attachment, relocation disruption, and assembly surplus, creating a windfall for private transferees at displaced homeowners' expense.

\textbf{Unworkable pretext test.} Lower courts have struggled to apply Kennedy's fact-intensive pretext inquiry because private profit and public tax-base gains are intertwined in nearly every economic-development project.

\section{Key Quotations}

\begin{quote}
``Promoting economic development is a traditional and long accepted function of government. There is, moreover, no principled way of distinguishing economic development from the other public purposes that we have recognized. \dots\ We emphasize that nothing in our opinion precludes any State from placing further restrictions on its exercise of the takings power.''
\end{quote}
\noindent --- Justice Stevens, Opinion of the Court, 545 U.S.\ at 484, 489.

\begin{quote}
``A court applying rational-basis review under the Public Use Clause should strike down a taking that, by a clear showing, is intended to favor a particular private party, with only incidental or pretextual public benefits.''
\end{quote}
\noindent --- Justice Kennedy, concurring, 545 U.S.\ at 491.

\begin{quote}
``The specter of condemnation hangs over all property. Nothing is to prevent the State from replacing any Motel 6 with a Ritz-Carlton, any home with a shopping mall, or any farm with a factory. \dots\ The beneficiaries are likely to be those citizens with disproportionate influence and power in the political process, including large corporations and development firms.''
\end{quote}
\noindent --- Justice O'Connor, dissenting, 545 U.S.\ at 503, 505.

\end{document}
